% id: 0489
% book: 라이트쎈 공통수학2 · 04 도형의 이동
% section: 실전 Training
% figure: none
% answer: ④
% answer_source: 해설
% variant_level: 0 (원본 전사)
% 이 파일은 build.mjs 가 items.json 에서 만든다. 고칠 때는 items.json 을 고친다.
\smash{\raisebox{1.5mm}{\dmkichul{라이트쎈 공통수학2 0489번 · 75쪽 · 교육청 기출}}}
좌표평면 위의 점~$\pt{A}(-3,\mkern3mu 4)$를 직선~$y=x$에 대하여 대칭이동한 점을 $\pt{B}$라 하고, 점~$\pt{B}$를 $x$축의 방향으로 $2$만큼, $y$축의 방향으로 $k$만큼 평행이동한 점을 $\pt{C}$라 하자. 세 점~$\pt{A}$, $\pt{B}$, $\pt{C}$가 한 직선 위에 있을 때, 실수 $k$의 값은?
\begin{choices32}{$-5$}{$-4$}{$-3$}{$-2$}{$-1$}\end{choices32}
